\documentclass[11pt]{article}
\usepackage[T1]{fontenc}
\usepackage{lmodern}
\usepackage[margin=1in]{geometry}
\usepackage{amsmath,amssymb,amsthm}
\usepackage{microtype}
\usepackage[hidelinks]{hyperref}
\newtheorem{theorem}{Theorem}
\newtheorem{proposition}{Proposition}
\newcommand{\Ric}{\operatorname{Ric}}
\newcommand{\Hess}{\operatorname{Hess}}
\newcommand{\tr}{\operatorname{tr}}
\newcommand{\dV}{\,dV}
\newcommand{\dA}{\,dA}
\title{A regular zero-area singularity with a nonzero weak Bianchi defect}
\author{Alec Kriebel\protect\\[3pt]
\small\href{https://orcid.org/0009-0001-9320-500X}{ORCID: 0009-0001-9320-500X}}
\date{Version 1.0 --- 12 October 2026}
\begin{document}
\maketitle

\begin{abstract}
We apply a negative-mass member of the known Dadhich--Kar--Mukherjee--Visser
metric family to the regular-zero-area-singularity conjecture in Burtscher's
Oberwolfach report. Its natural static slices have harmonically regular
zero-area singularities of mass $-2a$, but the homogeneous weak second Bianchi
identity fails under the boundary-reaching vector-test convention of
Burtscher--Kiessling--Tahvildar-Zadeh. We prove absolute convergence of the full
scalar contraction for every such compactly supported smooth test and compute
its exact boundary functional; a normalized radial test gives $-8\pi$.
The metric and its Einstein tensor are classical, and the flux method is
already present in the cited weak-identity work. The contribution is this
explicit counterexample application and its convergence check. The example
violates the null energy condition. The conclusion concerns the geometric
implication in the cited static-foliation framework, with no specified matter
model; it asserts neither regularity on every arbitrary tilted slice nor an
identity for tests smooth after collapsing the boundary sphere to a point.
\end{abstract}

\section{The geometric question and the test convention}

Burtscher's contribution to Oberwolfach Report 40/2021 describes admissible
timelike singularities using regular zero-area singularities (ZAS) of negative
mass on spacelike slices, and conjectures the homogeneous weak Einstein-tensor
identity for that class \cite[pp.~2232--2233]{OWR}. The cited technical framework
of Burtscher--Kiessling--Tahvildar-Zadeh (BKTZ) uses static slices and a retained
inner boundary \cite[Sections~1.3--1.4]{BKTZ}. This note addresses that geometric
implication with their boundary-reaching vector tests. The short report does
not give a fully formal definition for every nonsymmetric spacetime or
explicitly require regularity on every tilted hypersurface. We state the
foliation interpretation explicitly rather than claim to settle every
possible strengthening of its wording. The report's broader matter and
evolution program remains a separate question.

Fix $a>0$ and define
\[
 \overline M=\mathbb R_t\times\{x\in\mathbb R^3:|x|\ge a\},\qquad
 M=\operatorname{int}\overline M,\qquad
 \rho=|x|,\quad S_a=\{\rho=a\},\quad \phi=1-\frac a\rho.
\]
The smooth structure of $\overline M$ retains the sphere $S_a$. Our test space
is $\mathcal T=C_c^\infty(\overline M;T\overline M)$: smooth vector-field sections
with compact support in the manifold-with-boundary topology. Their normal
boundary component need not vanish. This is the convention in BKTZ,
Definition~2.1, the boundary-support footnote to Theorem~1.1, and the radial term
in equation~(2.5) \cite{BKTZ}. A section of the tangent bundle at a boundary is
not required to be tangent to that boundary.

We use signature $(-+++)$, $c=G_{\rm Newton}=1$, and
$G_{\mu\nu}=\Ric(g)_{\mu\nu}-\tfrac12 R_g g_{\mu\nu}$, with the Ricci convention
\[
R_{\mu\nu}=\partial_\lambda\Gamma^\lambda_{\mu\nu}
-\partial_\nu\Gamma^\lambda_{\mu\lambda}
+\Gamma^\lambda_{\mu\nu}\Gamma^\kappa_{\lambda\kappa}
-\Gamma^\kappa_{\mu\lambda}\Gamma^\lambda_{\nu\kappa}.
\]
For a metric smooth on $M$, write
\begin{equation}\label{eq:I}
 I_g(\psi)=\int_M G^\mu{}_{\nu}\nabla_\mu\psi^\nu\dV_g,
 \qquad \psi\in\mathcal T,
\end{equation}
whenever the full scalar contraction is absolutely integrable. On $M$ this
equals $\int_M G_{\mu\nu}\nabla^\mu\psi^\nu\dV_g$ pointwise. Lowering a test vector
with a singular metric need not preserve boundary smoothness; independently
smooth covectors are not substituted for $\mathcal T$. The homogeneous weak
identity is $I_g(\psi)=0$ for every test in this specified class.

\section{A classical metric and an explicit counterexample}

The metric used below and the affine-lapse family used in its calculation are
known. Dadhich--Kar--Mukherjee--Visser (DKMV) give
\cite[equations~(9)--(15)]{DKMV}
\[
 ds^2=-\left(\kappa+\lambda\frac{1-m/(2\bar r)}{1+m/(2\bar r)}\right)^2dt^2
       +\left(1+\frac{m}{2\bar r}\right)^4
        (d\bar r^2+\bar r^2d\Omega^2),
\]
along with its Einstein-tensor components. The substitution
\begin{equation}\label{eq:dictionary}
 (m,\bar r,\kappa,\lambda)=(-2a,\rho,A,B),\qquad
 N_v=\frac{\rho+a}{\rho-a},\qquad N=A+B N_v
\end{equation}
gives $g_N=-N^2dt^2+\phi^4\delta$, where $\delta$ is the Euclidean metric.
We take $A,B\ge0$, $A+B>0$, so $N>0$. DKMV's subsequent wormhole parameter
classification takes $m>0$; the negative-mass substitution here is an
application of their metric formula, not an attribution of this later ZAS or
weak-Bianchi conclusion to that classification. No new metric family or
Einstein tensor is claimed. The spatial negative-mass Schwarzschild ZAS is
also the standard example in Bray--Jauregui \cite{BJ}.

\begin{theorem}\label{thm:main}
For $a>0$, the metric
\begin{equation}\label{eq:mainmetric}
 g=-\phi^{-2}dt^2+\phi^4(d\rho^2+\rho^2d\Omega^2)
\end{equation}
is static, spherically symmetric and asymptotically flat on $M$. Every slice
of its natural static foliation is totally geodesic and has a harmonically
regular ZAS at $S_a$ of regular mass $-2a$. Its inner end is a genuine
curvature singularity, timelike in the limiting radial conformal sense.
For every $\psi\in\mathcal T$, the scalar in \eqref{eq:I} is absolutely
integrable and
\begin{equation}\label{eq:mainflux}
 I_g(\psi)=-\frac{2}{a^2}
 \int_{\mathbb R\times S_a}
       \langle\psi_{\rm spatial},n\rangle_\delta\,dt\dA_\delta,
 \qquad n=\frac{x}{\rho}.
\end{equation}
In particular there is a test in $\mathcal T$ with $I_g(\psi)=-8\pi$.
The effective tensor $T=G/(8\pi)$ violates the null energy condition.
\end{theorem}

The theorem disproves the proposed implication from regular negative-mass
ZAS slice geometry to the homogeneous identity under the stated foliation
and test convention. A specified matter system or an added energy condition
is not among its hypotheses. We prove the theorem through the following
geometric and analytic calculations; no numerical verification is used as a
proof of a quantified assertion.

\subsection{Regular ZAS and the timelike singular end}

The resolution $(\delta,\phi)$ extends smoothly to $S_a$, where $\phi$ vanishes
simply. The Euclidean unit normal into the exterior is $+\partial_\rho$, and
\[
 \partial_\rho\phi|_{S_a}=\frac1a>0,\qquad \Delta_\delta\phi=0.
\]
These are the regular and harmonic resolution conditions of
Bray--Jauregui, Definitions~6--7 \cite{BJ}. Their ZAS area condition concerns
approaching boundary graphs converging in $C^1$ (Definition~4). For every such
family the Euclidean areas stay bounded, while the supremum of $\phi^4$
on the surfaces tends to zero. Since $dA_h=\phi^4dA_\delta$ for
$h=\phi^4\delta$, their physical areas tend to zero. This checks the full
approaching-surface condition, not just coordinate spheres. Definition~10
then gives
\begin{equation}\label{eq:mass}
 m_{\rm reg}=-\frac14\left(\frac1\pi\int_{S_a}
       (\partial_\rho\phi)^{4/3}\dA_\delta\right)^{3/2}
 =-\frac14(4a^{2/3})^{3/2}=-2a.
\end{equation}
The time-independent spatial metric and zero shift give second fundamental
form $K=0$. These statements concern the natural static foliation.

The area radius satisfies
\begin{equation}\label{eq:areal}
 r=\rho\phi^2=\frac{(\rho-a)^2}{\rho},\qquad
 f=1+\frac{4a}{r}=N_v^2,
 \qquad h=f^{-1}dr^2+r^2d\Omega^2.
\end{equation}
Here $dr/d\rho=1-a^2/\rho^2>0$ on $M$; the change of variable is used only
on the smooth interior. At infinity $\phi^4=1-4a/\rho+O(\rho^{-2})$ and
$\phi^{-1}=1+a/\rho+O(\rho^{-2})$, with the corresponding differentiated
decay. Thus \eqref{eq:mainmetric} is asymptotically flat.

For its radial causal geometry put
\[
 s(\rho)=\int_a^\rho\phi(u)^3\,du
 \sim\frac{(\rho-a)^4}{4a^3}.
\]
The radial metric is $\phi^{-2}(-dt^2+ds^2)$, so its limiting conformal
boundary $s=0$ is timelike. Each nearby constant-$\rho$ tube has spacelike
normal and Lorentzian induced metric. A radial null geodesic with conserved
energy $E=\phi^{-2}dt/d\lambda>0$ obeys $|d\rho/d\lambda|=E/\phi$; its remaining
affine length is
\[
 \frac1E\int_a^\rho\phi(u)\,du
 =\frac{\rho-a-a\log(\rho/a)}{E}
 \sim\frac{(\rho-a)^2}{2aE}.
\]
The invariant computed below diverges at this finite-affine-length end.
The timelike statement is about this limiting causal structure; no
nondegenerate smooth extension of the physical four-metric to the cylinder
is asserted.

\subsection{Curvature calculation}

For completeness we recover the old tensor directly. For
$h=f^{-1}dr^2+r^2d\Omega^2$, its spatial Ricci and radial-Hessian mixed
eigenvalues are
\begin{align*}
 (\Ric(h)^r{}_r,\Ric(h)^\theta{}_\theta,\Ric(h)^\varphi{}_\varphi)
   &=\left(\frac{4a}{r^3},-\frac{2a}{r^3},-\frac{2a}{r^3}\right),\\
 (\Hess_h u)^r{}_r&=fu''+\frac{f'}2u',\qquad
 (\Hess_h u)^\theta{}_\theta=(\Hess_h u)^\varphi{}_\varphi=\frac{fu'}r.
\end{align*}
These follow from
$\Gamma^r_{rr}=-f'/(2f)$, $\Gamma^r_{\theta\theta}=-fr$,
$\Gamma^r_{\varphi\varphi}=-fr\sin^2\theta$,
$\Gamma^\theta_{r\theta}=\Gamma^\varphi_{r\varphi}=1/r$, and the usual
unit-sphere connection. With $q=N_v=\sqrt f$,
\[
 q'=-\frac{2a}{qr^2},\qquad
 q''=\frac{4a}{qr^3}-\frac{4a^2}{q^3r^4},
 \qquad \Hess_hq=q\Ric(h),\quad \Delta_hq=0.
\]
For $g_N=-N^2dt^2+h$, the additional time connection gives
\[
 \Ric(g_N)_{tt}=N\Delta_hN,\quad \Ric(g_N)_{ti}=0,\quad
 \Ric(g_N)_{ij}=\Ric(h)_{ij}-N^{-1}(\Hess_hN)_{ij}.
\]
Using $N=A+Bq$ and $R_h=0$, we obtain $R_{g_N}=0$, all off-diagonal
mixed Einstein entries zero, and
\begin{equation}\label{eq:familyG}
 G^t{}_t=0,\qquad
 (G^r{}_r,G^\theta{}_\theta,G^\varphi{}_\varphi)
 =\frac A N\left(\frac{4a}{r^3},-\frac{2a}{r^3},-\frac{2a}{r^3}\right).
\end{equation}
These agree with DKMV's equations~(9)--(11) under \eqref{eq:dictionary}.
The main lapse is $N=(1+q)/2=\phi^{-1}$, so
\begin{equation}\label{eq:mainG}
 (G^t{}_t,G^\rho{}_\rho,G^\theta{}_\theta,G^\varphi{}_\varphi)
 =\left(0,\frac{2a}{\rho^3\phi^5},
            -\frac{a}{\rho^3\phi^5},-\frac{a}{\rho^3\phi^5}\right).
\end{equation}
Since $R_g=0$, squaring the Ricci tensor yields
\begin{equation}\label{eq:invariant}
 \Ric(g)_{\mu\nu}\Ric(g)^{\mu\nu}
       =\frac{6a^2}{\rho^6\phi^{10}}\longrightarrow+\infty
       \quad(\rho\downarrow a).
\end{equation}
This is invariant curvature blow-up, despite the zero scalar curvature.
In a static orthonormal frame the energy density is zero and each angular
pressure is $-a/(8\pi\rho^3\phi^5)<0$. An angular null vector with frame
components $(1,0,1,0)$ therefore has $T(k,k)<0$, proving the stated NEC
violation.

\section{Absolute convergence and the exact defect}

We establish the full boundary functional for the classical family; its
main member gives \eqref{eq:mainflux}. BKTZ already use a limiting
boundary-flux calculation \cite[Definition~2.1 and the calculation following
equation~(2.5)]{BKTZ}. Here the necessary convergence is proved directly,
without invoking their sufficient theorem outside its hypotheses.

\begin{proposition}\label{prop:flux}
For $a>0$, $A,B\ge0$, $A+B>0$, let $g_N$ be defined by
\eqref{eq:dictionary}. For every $\psi\in\mathcal T$ the full contraction is
absolutely integrable and
\begin{equation}\label{eq:familyflux}
 I_{g_N}(\psi)=-\frac{4A}{a^2}
 \int_{\mathbb R\times S_a}
       \langle\psi_{\rm spatial},n\rangle_\delta\,dt\dA_\delta.
\end{equation}
Thus the homogeneous identity for every test in $\mathcal T$ holds within
this family exactly when $A=0$.
\end{proposition}

\begin{proof}
In exterior Cartesian coordinates $dV_{g_N}=N\phi^6\,dt\,d^3x$.
Equation~\eqref{eq:familyG} and $r^3=\rho^3\phi^6$ give the weighted mixed
spatial tensor
\begin{equation}\label{eq:D}
 D^i{}_j:=N\phi^6G^i{}_j
    =\frac{Aa}{\rho^3}(6n_i n_j-2\delta_{ij}).
\end{equation}
It is Euclidean symmetric, traceless, and smooth up to $\rho=a>0$.
Writing $\phi_i=\partial_i\phi$, the spatial connection is
\[
 \Gamma^j_{ik}=\frac2\phi
    (\delta^j_i\phi_k+\delta^j_k\phi_i-\delta_{ik}\phi^j),
 \qquad
 D^i{}_j\Gamma^j_{ik}
 =\frac2\phi\bigl((\tr D)\phi_k+D^i{}_k\phi_i-D^k{}_j\phi_j\bigr)=0.
\]
The time row and column of $G^\mu{}_{\nu}$ vanish, and
$\Gamma^j_{it}=0$. Hence, pointwise on $M$,
\begin{equation}\label{eq:scalar}
 G^\mu{}_{\nu}\nabla_\mu\psi^\nu\dV_{g_N}
       =D^i{}_j\partial_i\psi^j\,dt\,d^3x.
\end{equation}
On the compact support of any $\psi\in\mathcal T$, $D$ and the Cartesian
first derivatives of $\psi$ are bounded, including at the boundary.
The support has finite Euclidean measure. Thus the absolute value of the
right side is integrable. This proves convergence for arbitrary time and
angular dependence and avoids spherical-coordinate pole issues.

The radial and tangential eigenvalues of $D$ are
$D_r=4Aa/\rho^3$ and $D_T=-2Aa/\rho^3$. Consequently
\[
 \partial_iD^i{}_j
 =\left(D_r'+\frac2\rho(D_r-D_T)\right)n_j=0.
\]
Integrate \eqref{eq:scalar} by parts over $a+\varepsilon<\rho<R$, with $R$
and the time endpoints outside the support. The inner outward Euclidean
normal of this exterior domain is $-n$. The only boundary contribution is
\[
 -\frac{4Aa}{(a+\varepsilon)^3}
  \int_{\mathbb R\times S_{a+\varepsilon}}
       \langle\psi_{\rm spatial},n\rangle_\delta\,dt\dA_\delta.
\]
Absolute convergence of the interior integral and smoothness of the
boundary trace permit $\varepsilon\downarrow0$, proving
\eqref{eq:familyflux}. At $A=0$ the Einstein tensor is identically zero.
For $A>0$ the radial test below has nonzero functional, proving the final
assertion.
\end{proof}

Choose $\eta\in C_c^\infty(\mathbb R)$ with $\int\eta\,dt=1$ and
$\chi\in C_c^\infty([a,\infty))$ equal to $1$ near $a$. Set
\begin{equation}\label{eq:test}
 \psi=\eta(t)\chi(\rho)\partial_\rho.
\end{equation}
Since $\partial_\rho=x^i\partial_{x^i}/\rho$ and $a>0$, this is smooth on
$\overline M$ and has a nonzero normal boundary trace. Its spherical
components are angular independent, meeting also the convenient restriction
used in BKTZ's equation~(2.5). Formula~\eqref{eq:familyflux} gives
$I_{g_N}(\psi)=-16\pi A$, hence $-8\pi$ at $A=B=1/2$.
As a direct sign and density check, its complete scalar density is
\[
 4Aa\eta(t)\left(\frac{\chi'(\rho)}\rho-
                     \frac{\chi(\rho)}{\rho^2}\right)
        \,dt\,d\rho\,d\Omega
 =4Aa\eta(t)\partial_\rho(\chi/\rho)\,dt\,d\rho\,d\Omega.
\]
The boundary endpoint is $\chi(a)/a=1/a$, which gives the same negative
sign. This completes the proof of Theorem~\ref{thm:main}.

Individual connection summands in this radial computation can each have
logarithmically divergent integrals. Equation~\eqref{eq:scalar} combines
the ordinary smooth summands on $M$ before integration. It does not
separately integrate divergent terms, take a principal value, multiply
tensor distributions, or assign a nonsingular boundary value to the
physical metric.

\section{Scope, checks, and disclosure}

The counterexample has nonzero order $(\rho-a)^{-5}$ mixed Einstein
components. In particular, with $x=\rho-a$,
$x^5G^\rho{}_\rho\to2a^3\ne0$. BKTZ Theorem~1.1 assumes the strictly
improved bound $O(x^{-5+\kappa})$ for some $\kappa>0$ \cite{BKTZ}; it does
not cover this borderline case. The old boundary-flux method and the
independent convergence proof are distinct from an invocation of that
sufficient theorem. The vacuum control $A=0,B>0$ is the negative-mass
Schwarzschild vacuum spacetime after time rescaling and has zero functional.
Sharing its spatial metric does not make the main lapse vacuum.

The static-slice condition is essential to the claim's formulation.
Bray--Jauregui use a time-symmetric setting and discuss extensions beyond
it separately \cite[p.~3 and Section~6.6 of the arXiv version]{BJ}. For
example, a uniformly spacelike local tilt $t=t_0+q s(\rho)$, $|q|<1$, has
induced metric
\[
 h_q=\phi^4\bigl((1-q^2)d\rho^2+\rho^2d\Omega^2\bigr),
 \qquad m_{\rm reg}(h_q)=-\frac{2a}{1-q^2}<0.
\]
Its resolution remains smooth and positive. Arbitrary nearly null tilts
need not have that property: $t=u(\rho)$ with
$u'=\phi^3\sqrt{1-\phi^2}$ is interior-spacelike but has induced metric
$\phi^6d\rho^2+\phi^4\rho^2d\Omega^2$. Any simple defining conformal factor
in the retained cylinder structure leaves a zero radial eigenvalue in the
resolved boundary metric. No all-arbitrary-tilted-slices assertion is made.

Tests confined to the open interior give the ordinary classical identity,
and tests smooth after collapsing $S_a$ to a point have a different
smoothness requirement. In particular the radial test \eqref{eq:test}
must not be replaced through that collapse. The effective matter violates
NEC, so neither an energy-constrained version nor a specified electromagnetic
matter system is resolved. The conclusion is the failure of the geometric
implication in the stated boundary/static framework.

The accompanying verification materials bind the thirteen available
historical source files and preserve their original bytes. Independent
analytic derivations and exact symbolic or rational-jet calculations check
the tensor, mass, scalar cancellation, boundary sign and relevant controls.
Historical 17-check and 35-check outputs reproduce in ordinary and optimized
Python. The author's historical final assertion alone disappears in
optimized mode; a separate explicit replay guard checks every Boolean,
the expected counts, success flags and exact output bytes in both modes.
Eleven extant historical payload references are checked; one standalone
historical author-manifest hash has no supplied manifest path or bytes and
is not claimed verified. These provenance facts supplement the analytic
proof rather than replace it.

A bounded independent literature search found no earlier printed application
of this metric resolving the exact report conjecture in the inspected
routes. This is not a certificate of historical first priority or continued
openness. The metric, tensor, spatial ZAS example and boundary-flux method
are attributed to the sources above. The supporting priority ledger records
search coverage and its limits.

\paragraph{AI use and review status.}
AI tools were used extensively in developing, writing, reproducing and
adversarially checking this work, including parallel AI reviews. This is an
unrefereed research note. No independent human peer review or formal
proof-assistant verification is claimed.

\begin{thebibliography}{9}
\bibitem{OWR}
A. Burtscher, \emph{Matter particles, naked singularities, and a weak second
Bianchi identity}, in \emph{Mathematical Aspects of General Relativity},
Oberwolfach Reports \textbf{18} (2021), no.~3, Report~40, pp.~2231--2233.
\href{https://doi.org/10.4171/OWR/2021/40}{doi:10.4171/OWR/2021/40}.
The volume/workshop year is 2021; the publisher lists publication on
26 November 2022.

\bibitem{BKTZ}
A. Burtscher, M. K.-H. Kiessling and A. S. Tahvildar-Zadeh,
\emph{Weak second Bianchi identity for static, spherically symmetric
spacetimes with timelike singularities}, Classical and Quantum Gravity
\textbf{38} (2021), 185001.
\href{https://doi.org/10.1088/1361-6382/ac1853}{doi:10.1088/1361-6382/ac1853};
\href{https://arxiv.org/abs/1901.00813v3}{arXiv:1901.00813v3}.
Numbered locations in this note refer to that arXiv version.

\bibitem{DKMV}
N. Dadhich, S. Kar, S. Mukherjee and M. Visser,
\emph{$R=0$ spacetimes and self-dual Lorentzian wormholes},
Physical Review D \textbf{65} (2002), 064004.
\href{https://doi.org/10.1103/PhysRevD.65.064004}{doi:10.1103/PhysRevD.65.064004};
\href{https://arxiv.org/abs/gr-qc/0109069}{arXiv:gr-qc/0109069} (2001).
Published 7 February 2002. Equation numbers refer to the original arXiv
manuscript.

\bibitem{BJ}
H. L. Bray and J. L. Jauregui, \emph{A geometric theory of zero area
singularities in general relativity}, Asian Journal of Mathematics
\textbf{17} (2013), 525--559;
\href{https://arxiv.org/abs/0909.0522v2}{arXiv:0909.0522v2}.
Definition and section locations refer to that arXiv version.
\end{thebibliography}
\end{document}
